%% kapa: exam=1 round=24 year=2013 period=1 form=A subject=civil_law subject_ko=민법 seq=19 number=19
%% source: Q-Net 감정평가사 기출문제 (cst00309) artlSeq=5073802 file=extracted/1008/5073802/2013년 제24회 감정평가사 1차시험 기출문제/1교시(A형).hwp
%% conversion: kordoc 4.24.51 -> tools/gen1.py (auto); review=proofread-sonnet; flags=none
\begin{question}{19}
소멸시효의 중단에 관한 설명으로 옳지 않은 것은? (다툼이 있으면 판례에 의함)
\begin{choices}
  \item 시효기간만료 전에 응소함으로써 시효가 중단되는 경우 소멸시효기간이 만료된 후라도 사실심 변론종결 전에는 언제든지 시효중단을 주장할 수 있다.
  \item 시효중단의 효력은 당사자 및 그 승계인 사이에서만 효력이 있다.
  \item 시효중단의 효력 있는 승인에는 상대방의 권리에 관한 처분의 능력이나 권한 있음을 요한다.
  \item 사망한 사람을 피신청인으로 한 가압류결정은 상속인에 대하여 시효중단의 효력이 없다.
  \item 가처분은 시효의 이익을 받은 자에 대하여 하지 아니한 때에는 이를 그에게 통지한 후가 아니면 시효중단의 효력이 없다.
\end{choices}
\end{question}
\answer{3}
